% Author: Maximiliano Ochoa Trancoso (Colegio de México)
% Modified template from Aaron Wolf (Northwestern University)
\documentclass[aspectratio=169]{beamer}
\makeatletter
\def\input@path{{theme/}}
\makeatother

\title{Metodología de la Investigación}
\date{Agosto 2026}
\author{Maximiliano Ochoa Trancoso}

% Theme options:
% \usetheme[
%           pattern=triangle,
%           color=wcprimary,
%           font=geo,
%           agendastyle=pattern,
%           breadcrumb=on
%   ]{colmex}
% pattern: facet (default), facetrand, triangle , triangleiso, trianglemesh,
%          squares, circles, hexraise, hexrand, network, quilt, plain 
% There are some issues with the compilation time of some patterns
% color: any named color 
% font: noto (Noto Sans), Geomanist (default) or installed (system fonts)
% agendastyle: plain  or pattern (default)
% breadcrumb: on or none (default)
\usetheme{colmex}

    \titlegraphic{
        \makebox[\textwidth]{
            \includegraphics[height=1.7cm]{logos/logoColmex.png}
            \hfill \hspace{0.5cm}
            \includegraphics[height=1.8cm]{logos/logoCEE.png}
        }
    }

\begin{document}

    \begin{frame}
    \titlepage
    \end{frame}
    


    \begin{frame}[plain]
        \tableofcontents
    \end{frame}

\section{Blocks}

    \begin{frame}{Blocks}
        
        \begin{block}{Primary Block}
            The standard Beamer \texttt{block} environment uses the material-style outlined box with a floating title label.
        \end{block}
        \begin{alertblock}{Alert Block}
            \texttt{alertblock} uses the theme alert color.
        \end{alertblock}
        \begin{exampleblock}{Example Block}
            \texttt{exampleblock} uses the theme example color.
        \end{exampleblock}
    
    \end{frame}
    
    \begin{frame}{Blocks}
        \begin{palertblock}{Pattern Alert Block}
            Patterned alert block.
        \end{palertblock}
        \begin{pexampleblock}{Pattern Example Block}
            Patterned example block.
        \end{pexampleblock}
        \begin{pblock}{Pattern Block}
        If you want to use the main background pattern inside a tcolorbox, the theme also includes \texttt{pblock}, \texttt{palertblock}, and \texttt{pexampleblock}.
        \end{pblock}
    \end{frame}
    
    \begin{frame}{Custom Block Colors}
        \begin{mblock}[colframe=HarvestOrange,coltitle=HarvestOrange]{Custom Block}
            This block uses a custom color.
        \end{mblock}
        \begin{pblock}[VintageLavender]{Custom Pattern Block}
            This patterned block uses a custom color.
        \end{pblock}
    \end{frame}

\section{Frames}
    
% Blank Frame
    \begin{frame}[plain]{}
        The \texttt{[plain]} option produces a fully blank frame with no header, footer, or background. Use it for title cards, full-bleed figures, or transition slides:
        \\ ~ \\
        \texttt{\textbackslash begin\{frame\}[plain]\{\}} \\
        \texttt{...} \\
        \texttt{\textbackslash end\{frame\}}
    \end{frame}

    
% Standout frame
    \begin{frame}{Standout Slides}
        The \texttt{\textbackslash standout\{\}} command creates a full-pattern slide with centered white text — useful for section breaks, closing slides, or emphasis.
        \\ ~ \\
        \scriptsize
        \texttt{\textbackslash standout\{Appendix\}}
        \normalsize
    \end{frame}
    
    \standout{Appendix}

\appendix




\end{document}
